%% Project 2, FYS-STK3155/4155, fall 2026 -- first iteration (2026-10-06)
%% Converted from the 2025 DocOnce source doc/BookML/Project2.do.txt.
%% Changes relative to 2025: links point to the public repository
%% https://github.com/EducationalMaterialUiO/MachineLearningUiO (weekly
%% material in doc/WeeklyMaterial/weekNN, lecture notes as jupyter-book);
%% the preamble (reports, references, LLM declaration) is the same as in
%% Project 1; a list of insights to take away; references to the relevant
%% sections of Chapter 8 (and Chapters 4 and 5) of the lecture notes; the
%% learning rate is called gamma as in the lecture notes and Project 1.
%% The notebook variant Project2.ipynb in this folder carries the same text
%% (generated from this file by build_project2_notebook.py).
%% Build:  pdflatex Project2.tex  (twice)
\documentclass[11pt,a4paper]{article}
\usepackage[T1]{fontenc}
\usepackage[utf8]{inputenc}
\usepackage{lmodern}
\usepackage{amsmath,amssymb,bm}
\usepackage[margin=2.5cm]{geometry}
\usepackage{parskip}
\usepackage{enumitem}
\usepackage{listings}
\usepackage{xcolor}
\usepackage[normalem]{ulem}   % \uline: underline that can break across lines
\usepackage[colorlinks=true,linkcolor=blue!60!black,urlcolor=blue!60!black]{hyperref}
\usepackage{xurl}             % let the long URLs break anywhere

\definecolor{codebg}{rgb}{0.97,0.97,0.97}
\lstset{language=Python,basicstyle=\ttfamily\small,backgroundcolor=\color{codebg},
  keywordstyle=\color{blue!70!black}\bfseries,commentstyle=\color{green!45!black},
  stringstyle=\color{red!60!black},showstringspaces=false,frame=single,
  framerule=0pt,xleftmargin=4pt,xrightmargin=4pt,breaklines=true,columns=fullflexible}

\title{Project 2 on Machine Learning\\[4pt]
  \large Classification and regression: writing our own neural network code\\[4pt]
  \normalsize Deadline November 6 (midnight), 2026}
\author{Data Analysis and Machine Learning FYS-STK3155/FYS4155\\
  Department of Physics, University of Oslo, Norway}
\date{October 6, 2026}

\begin{document}
\maketitle

\section*{Deliverables}

First, join a group in Canvas with your group partners.  Pick an available
group for Project 2 in the \textbf{People} page.

In Canvas, deliver as a group and include:
\begin{itemize}[nosep]
  \item A PDF of your report which follows the guidelines covered below, in
        the exercises of week 39 and in the Tuesday session of week 40.
        Additional requirements include:
  \begin{itemize}[nosep]
    \item It should be around 5000 words; use the word counter in Overleaf
          for this.  This often corresponds to 10--12 pages.  References and
          appendices are excluded from the word count.  Appendices should
          read like regular text and can include additional calculations you
          have made.
    \item It should include around 10--15 figures.  You can include more
          figures in appendices and/or as supplemental material in your
          repository.  Including additional figures in appendices requires
          also that you add text which explains the additional material.
  \end{itemize}
  \item A comment linking to your GitHub repository (or a folder in one of
        your GitHub repositories) for this project.  The repository must
        include
  \begin{itemize}[nosep]
    \item a PDF file of the report;
    \item a folder named \texttt{Code}, where you put Python files for your
          functions and notebooks for reproducing your results.  Remember to
          use a seed for generating random data and for train-test splits
          when generating final results;
    \item a \texttt{README} file with the names of the group members, a short
          description of the project, a description of how to install the
          required packages to run your code (from a
          \texttt{requirements.txt} file or similar, such as a plain text
          description), and the names and descriptions of the various
          notebooks in the \texttt{Code} folder and the results they produce.
  \end{itemize}
\end{itemize}

\section*{Preamble: note on writing reports, using reference material, AI and other tools}

We want you to answer the three different projects by handing in
reports written like a standard scientific/technical report.  The links at
\url{https://github.com/EducationalMaterialUiO/MachineLearningUiO/tree/main/doc/Projects}
contain more information.  There you can find examples of previous reports,
the projects themselves, how we grade reports etc.  How to write reports will
also be discussed during the various lab sessions.  Please do ask us if you
are in doubt. In that folder you will find how we evaluate the projects. You will receive a feedback with a final score for each project.

When using codes and material from other sources, you should refer to these
in the bibliography of your report, indicating wherefrom you for example got
the code, whether this is from the lecture notes, software like
\textbf{Scikit-Learn}, \textbf{JAX}, \textbf{TensorFlow}, \textbf{PyTorch} or other
sources.  These sources should always be cited correctly.  How to cite some of
the libraries is often indicated from their corresponding GitHub sites or
websites; see for example how to cite Scikit-Learn at
\url{https://scikit-learn.org/dev/about.html}.

You may use Large Language models (LLMs) like \href{https://openai.com/chatgpt/}{ChatGPT},
\href{https://claude.ai}{Claude} or similar as a \emph{supportive measure} when
writing the report and developing your code, that is, as a tool that helps you
understand, check and improve work which is your own, and not as a substitute
for doing the analysis, writing the code or drawing the conclusions yourself.
In the report, you should after the conclusions, write a declaration on how you have used LLMs following the instructions at \url{https://github.com/EducationalMaterialUiO/MachineLearningUiO/blob/main/LLM_Usage_Declaration_Guidelines.md}. Failing to do so, \uline{leads to an automatic deduction of five (5) points out of the final score of 100 for the  project}.
Remember that \emph{you} are responsible for every
line of code and every claim in the report: be prepared to explain and defend
all of it.
If you use LLMs, please take a look at the guidelines at \url{https://github.com/EducationalMaterialUiO/MachineLearningUiO/tree/main/doc/UsingLLMs} on how to use these tools efficiently.

If you would like to study other data sets, feel free to propose other sets.
What we have proposed here are mere suggestions from our side.  If you opt for
another data set, consider using a set which has been studied in the
scientific literature.  This makes it easier for you to compare and analyze
your results.  Comparing with existing results from the scientific literature
is also an essential element of the scientific discussion.  The University of
California at Irvine with its Machine Learning repository at
\url{https://archive.ics.uci.edu/ml/index.php} is an excellent site to look up
for examples and inspiration.  \href{https://www.kaggle.com/}{Kaggle.com} is
an equally interesting site.  Feel free to explore these sites and/or ask us
for inputs relevant to your discipline.

\section*{Classification and regression, writing our own neural network code}

The main aim of this project is to study both classification and regression
problems by developing our own feed-forward neural network (FFNN) code.  The
lecture material and the exercises of weeks 41 and 42 (see
\url{https://github.com/EducationalMaterialUiO/MachineLearningUiO/tree/main/doc/WeeklyMaterial/week41}
and
\url{https://github.com/EducationalMaterialUiO/MachineLearningUiO/tree/main/doc/WeeklyMaterial/week42}),
together with Chapter 8 of the lecture notes (see
\url{https://educationalmaterialuio.github.io/MachineLearningUiO/doc/LectureNotes/_build/html/chapter8.html}),
should contain enough information for you to get started with writing your
own code.

We will also reuse our codes on gradient descent methods from project 1.

The data sets that we propose here are (the default sets)
\begin{itemize}[nosep]
  \item \textbf{Regression} (fitting a continuous function).  In this part
        you will need to bring back your results from project 1 and compare
        these with what you get from the neural network code to be developed
        here.  The data sets could be
  \begin{itemize}[nosep]
    \item the simple one-dimensional Runge function from project 1, that is
          $f(x) = \frac{1}{1+25x^2}$.  We recommend using a simpler function
          when developing your neural network code for regression problems.
          Feel however free to discuss and study other functions, such as the
          two-dimensional Runge function
          $f(x,y)=\left[(10x - 5)^2 + (10y - 5)^2 + 1 \right]^{-1}$, or even
          more complicated two-dimensional functions (see the supplementary
          material of \url{https://www.nature.com/articles/s41467-025-61362-4}
          for an extensive list of two-dimensional functions).
  \end{itemize}
  \item \textbf{Classification.}  We will consider a multiclass
        classification problem given by the full MNIST data set of
        handwritten digits, see
        \url{https://www.tensorflow.org/datasets/catalog/mnist} for a
        description of the data set and how to cite it.
\end{itemize}

The project is organized so that it follows the lectures and exercises of
weeks 41 and 42 and Chapter 8 of the lecture notes:
\begin{itemize}[nosep]
  \item Part a): the cost functions, their gradients and the activation
        functions we need (Sections
        \href{https://educationalmaterialuio.github.io/MachineLearningUiO/doc/LectureNotes/_build/html/chapter8.html#activation-functions}{8.6}
        and
        \href{https://educationalmaterialuio.github.io/MachineLearningUiO/doc/LectureNotes/_build/html/chapter8.html#the-cost-function-and-the-optimisation-problem}{8.7},
        and Sections 5.3 and 5.6 of Chapter 5 for the cross entropy and the
        softmax function).
  \item Part b): the feed-forward pass and the back propagation algorithm,
        and the first comparison with OLS from project 1 (Sections
        \href{https://educationalmaterialuio.github.io/MachineLearningUiO/doc/LectureNotes/_build/html/chapter8.html#notation-and-the-feed-forward-pass}{8.5},
        \href{https://educationalmaterialuio.github.io/MachineLearningUiO/doc/LectureNotes/_build/html/chapter8.html#backpropagation}{8.8},
        \href{https://educationalmaterialuio.github.io/MachineLearningUiO/doc/LectureNotes/_build/html/chapter8.html#weight-initialisation}{8.10}
        and
        \href{https://educationalmaterialuio.github.io/MachineLearningUiO/doc/LectureNotes/_build/html/chapter8.html#a-neural-network-from-scratch}{8.13};
        weeks 40 and 41).
  \item Parts c) and d): activation functions, depth and width, and the
        $L_1$ and $L_2$ penalties (Sections
        \href{https://educationalmaterialuio.github.io/MachineLearningUiO/doc/LectureNotes/_build/html/chapter8.html#activation-functions}{8.6},
        \href{https://educationalmaterialuio.github.io/MachineLearningUiO/doc/LectureNotes/_build/html/chapter8.html#vanishing-and-exploding-gradients}{8.9}
        and
        \href{https://educationalmaterialuio.github.io/MachineLearningUiO/doc/LectureNotes/_build/html/chapter8.html#regularisation-and-hyperparameters}{8.12};
        week 42).
  \item Part e): classification with the softmax output layer and the
        cross entropy (Sections
        \href{https://educationalmaterialuio.github.io/MachineLearningUiO/doc/LectureNotes/_build/html/chapter8.html#classification-with-neural-networks}{8.11}
        and
        \href{https://educationalmaterialuio.github.io/MachineLearningUiO/doc/LectureNotes/_build/html/chapter8.html#examples}{8.14},
        and Section 5.12 on how to measure the quality of a classifier;
        weeks 41 and 42).
  \item Part f): the critical evaluation (Section
        \href{https://educationalmaterialuio.github.io/MachineLearningUiO/doc/LectureNotes/_build/html/chapter8.html#limitations-and-a-top-down-view}{8.15}).
\end{itemize}
The analytical parts (part a) belong in the theory section of your report.
The project-style exercise at the end of Chapter 8 (Section
\href{https://educationalmaterialuio.github.io/MachineLearningUiO/doc/LectureNotes/_build/html/chapter8.html#exercises}{8.18})
follows the same steps as this project and can be used as a checklist.

\section*{What you should take away from this project}

Beyond the figures and tables in your report, these are the insights we
want you to be able to state, explain and defend after working through
the project.  Use them as a checklist when you write your conclusions.
\begin{enumerate}[nosep]
  \item \textbf{A neural network replaces hand-made features by learned
        ones.}  In project 1 the design matrix contained powers of $x$ that
        you chose.  A network takes the raw input and learns its own
        features in the hidden layers; the universal approximation theorem
        (Section
        \href{https://educationalmaterialuio.github.io/MachineLearningUiO/doc/LectureNotes/_build/html/chapter8.html#the-universal-approximation-theorem}{8.4})
        tells you that this is possible, not that it is easy.
  \item \textbf{Backpropagation is reverse-mode automatic differentiation
        applied to a layered function.}  The four equations of Section
        \href{https://educationalmaterialuio.github.io/MachineLearningUiO/doc/LectureNotes/_build/html/chapter8.html#the-four-equations-of-backpropagation}{8.8.4}
        are nothing but the chain rule organized so that one forward and one
        backward sweep give the gradient with respect to every weight and
        bias, at a cost of roughly two forward passes, independently of the
        number of parameters.  You should be able to derive them and to check
        your implementation against \texttt{jax.grad}.
  \item \textbf{The output activation and the cost function come in matched
        pairs.}  Linear output with the MSE, sigmoid with the binary cross
        entropy and softmax with the multiclass cross entropy all give the
        same output error $\bm{\delta}^L=\bm{a}^L-\bm{y}$ (up to the
        normalization).  Mixing the pairs (for instance a sigmoid output with
        the MSE) gives tiny gradients exactly where the network is
        confidently wrong.
  \item \textbf{The optimizer does not care where the gradient comes from.}
        Plain gradient descent, stochastic gradient descent, RMSprop and Adam
        from project 1 carry over unchanged; what changes is that the cost
        function is no longer convex, so the learning rate, the
        initialization and the minibatch size now decide whether you converge
        at all, and to which minimum.
  \item \textbf{Initialization and scaling are not details.}  Unscaled
        inputs or weights of the wrong size saturate the sigmoid, and the
        gradient vanishes (Sections
        \href{https://educationalmaterialuio.github.io/MachineLearningUiO/doc/LectureNotes/_build/html/chapter8.html#vanishing-and-exploding-gradients}{8.9}
        and
        \href{https://educationalmaterialuio.github.io/MachineLearningUiO/doc/LectureNotes/_build/html/chapter8.html#weight-initialisation}{8.10}).
        Standardized inputs,
        Xavier or He initialization and the ReLU family exist precisely to
        keep the gradients alive.
  \item \textbf{Hyperparameters must be searched systematically, and on
        validation data.}  The learning rate $\gamma$, the penalty $\lambda$,
        the depth, the width and the activation function interact.  A grid
        of two parameters at a time, shown as a heat map, is far more
        informative than a handful of hand-picked runs, and the test set is
        used only once, at the end.
  \item \textbf{Overfitting looks the same as in project 1.}  A growing gap
        between training and test error as the number of parameters grows is
        the same bias-variance story as for the polynomial degree, and the
        $L_1$ and $L_2$ penalties play the same role as in Ridge and Lasso
        regression; early stopping is a form of regularization that costs
        nothing.
  \item \textbf{More flexible is not automatically better.}  For a smooth
        one-dimensional function, OLS with a well-chosen polynomial degree is
        a closed-form, cheap and hard-to-beat benchmark.  The neural network
        earns its cost only when the features are not known in advance, as
        for the pixels of MNIST.  Your report should say when each method is
        the right tool.
  \item \textbf{A number is meaningless without a baseline.}  An accuracy
        of 97\% on MNIST is good only if you know that guessing gives 10\%,
        that logistic regression on raw pixels gives roughly 92\% and that
        convolutional networks reach above 99\%.  A confusion matrix tells
        you which errors remain; the accuracy alone does not.
  \item \textbf{A code you cannot verify is a code you cannot trust.}
        Gradient checks against automatic differentiation, recovering the
        closed-form OLS solution with a linear network, and comparing with
        Scikit-Learn or PyTorch are the tests that turn your implementation
        into a result you can defend.  Seeds and documented hyperparameters
        make the result reproducible.
\end{enumerate}

We will start with a regression problem and we will reuse our codes on
gradient descent methods from project 1.

\subsection*{Part a): Analytical warm-up}

When using our gradient machinery from project 1, we will need the
expressions for the cost/loss functions and their respective gradients.  The
functions whose gradients we need are:
\begin{enumerate}[nosep]
  \item the mean-squared error (MSE) with and without the $L_1$ and $L_2$
        norms (regression problems);
  \item the binary cross entropy (aka log loss) for binary classification
        problems with and without the $L_1$ and $L_2$ norms;
  \item the multiclass cross entropy cost/loss function (aka softmax cross
        entropy or just softmax loss function).
\end{enumerate}
Set up these three cost/loss functions and their respective derivatives and
explain the various terms.  In this project you will however only use the
MSE and the softmax cross entropy.  The cross entropy and the softmax
function are discussed in Sections 5.3 and 5.6 of the lecture notes, the
cost functions for neural networks in Section
\href{https://educationalmaterialuio.github.io/MachineLearningUiO/doc/LectureNotes/_build/html/chapter8.html#the-cost-function-and-the-optimisation-problem}{8.7},
and the regularization terms in Section
\href{https://educationalmaterialuio.github.io/MachineLearningUiO/doc/LectureNotes/_build/html/chapter8.html#regularisation-and-hyperparameters}{8.12}.

We will test three activation functions for our neural network setup, these
are
\begin{enumerate}[nosep]
  \item the Sigmoid (aka \emph{logit}) function,
  \item the RELU function and
  \item the Leaky RELU function.
\end{enumerate}
Set up their expressions and their first derivatives.  The activation
functions are discussed in Section
\href{https://educationalmaterialuio.github.io/MachineLearningUiO/doc/LectureNotes/_build/html/chapter8.html#activation-functions}{8.6}
of the lecture notes (the sigmoid in Section 8.6.1, the RELU family in
Section 8.6.2), and in the lecture material of week 41 at
\url{https://github.com/EducationalMaterialUiO/MachineLearningUiO/tree/main/doc/WeeklyMaterial/week41}.

\subsection*{Reminder about the gradient machinery from project 1}

In the setup of a neural network code you will need your gradient descent
codes from project 1.  For neural networks we will recommend using
stochastic gradient descent (Section 4.7 of the lecture notes) with either
the RMSprop or the Adam algorithms for updating the learning rates (Sections
4.10 and 4.11).  But you should feel free to try plain gradient descent as
well.  As in project 1 and in the lecture notes we call the learning rate
$\gamma$ (many texts call it $\eta$).

We recommend reading chapter 8 on optimization from the textbook of
Goodfellow, Bengio and Courville at \url{https://www.deeplearningbook.org/}.
This chapter contains many useful insights and discussions on the
optimization part of machine learning.  A useful reference on the back
propagation algorithm is Nielsen's book at
\url{http://neuralnetworksanddeeplearning.com/}.

You will find the Python
\href{https://seaborn.pydata.org/generated/seaborn.heatmap.html}{Seaborn
package} useful when plotting the results as function of the learning rate
$\gamma$ and the hyperparameter $\lambda$.

\subsection*{Part b): Writing your own neural network code}

Your aim now, and this is the central part of this project, is to write your
own FFNN code implementing the back propagation algorithm discussed in the
lecture material of week 41 (see
\url{https://github.com/EducationalMaterialUiO/MachineLearningUiO/tree/main/doc/WeeklyMaterial/week41})
and week 42 (see
\url{https://github.com/EducationalMaterialUiO/MachineLearningUiO/tree/main/doc/WeeklyMaterial/week42}),
and in Section
\href{https://educationalmaterialuio.github.io/MachineLearningUiO/doc/LectureNotes/_build/html/chapter8.html#backpropagation}{8.8}
of the lecture notes.  The notation and the feed-forward pass are set up in
Section
\href{https://educationalmaterialuio.github.io/MachineLearningUiO/doc/LectureNotes/_build/html/chapter8.html#notation-and-the-feed-forward-pass}{8.5},
the general dense-layer case and the four equations of backpropagation in
Sections
\href{https://educationalmaterialuio.github.io/MachineLearningUiO/doc/LectureNotes/_build/html/chapter8.html#the-general-case-dense-layers}{8.8.3}
and
\href{https://educationalmaterialuio.github.io/MachineLearningUiO/doc/LectureNotes/_build/html/chapter8.html#the-four-equations-of-backpropagation}{8.8.4},
and the full algorithm in Section
\href{https://educationalmaterialuio.github.io/MachineLearningUiO/doc/LectureNotes/_build/html/chapter8.html#the-algorithm}{8.8.5}.
A complete network from scratch, including a check of the gradient against
automatic differentiation, is discussed in Section
\href{https://educationalmaterialuio.github.io/MachineLearningUiO/doc/LectureNotes/_build/html/chapter8.html#a-neural-network-from-scratch}{8.13}.
The exercises of week 41 (\texttt{exercisesweek41.ipynb}) give you a
template for the code: samples are stored in the rows of the input, the
weight matrix of a layer has shape (inputs, outputs), and the feed-forward
step of a layer reads \texttt{a @ W + b}.

We will focus on a regression problem first, using the one-dimensional Runge
function
\[
  f(x) = \frac{1}{1+25x^2},
\]
from project 1.

Use only the mean-squared error as cost function (no regularization terms)
and write an FFNN code for a regression problem with a flexible number of
hidden layers and nodes using only the Sigmoid function as activation
function for the hidden layers.  Initialize the weights using a normal
distribution (see Section
\href{https://educationalmaterialuio.github.io/MachineLearningUiO/doc/LectureNotes/_build/html/chapter8.html#weight-initialisation}{8.10}
for how the width of that distribution should depend on the number of inputs
of a layer).  How would you initialize the biases?  And which activation
function would you select for the final output layer?  And how would you set
up your design/feature matrix?  Hint: does it have to represent a polynomial
approximation as you did in project 1?

Train your network and compare the results with those from your OLS
regression code from project 1 using the one-dimensional Runge function.
When comparing your neural network code with the OLS results from project 1,
use the same data sets which gave you the best MSE score.  Moreover, use the
polynomial order from project 1 that gave you the best result.  Compare these
results with your neural network with one and two hidden layers using $50$
and $100$ hidden nodes, respectively.

Comment your results and give a critical discussion of the results obtained
with the OLS code from project 1 and your own neural network code.  Make an
analysis of the learning rates employed to find the optimal MSE score.  Test
both stochastic gradient descent with RMSprop and Adam and plain gradient
descent with different learning rates.

You should, as you did in project 1, scale your data (Section 3.13 of the
lecture notes).

\subsection*{Part c): Testing different activation functions and depths of the neural network}

You should also test different activation functions for the hidden layers.
Try out the Sigmoid, the RELU and the Leaky RELU functions and discuss your
results (Section
\href{https://educationalmaterialuio.github.io/MachineLearningUiO/doc/LectureNotes/_build/html/chapter8.html#which-activation-should-we-use}{8.6.4}
of the lecture notes discusses which activation to use, and Section
\href{https://educationalmaterialuio.github.io/MachineLearningUiO/doc/LectureNotes/_build/html/chapter8.html#vanishing-and-exploding-gradients}{8.9}
why the sigmoid makes the gradients vanish in deep networks).  Test your
results as functions of the number of hidden layers and nodes.  Do you see
signs of overfitting?  It is optional in this project to perform a
bias-variance trade-off analysis.

\subsection*{Part d): Testing different norms}

Finally, still using the one-dimensional Runge function, add now the
hyperparameters $\lambda$ with the $L_2$ and $L_1$ norms (Section
\href{https://educationalmaterialuio.github.io/MachineLearningUiO/doc/LectureNotes/_build/html/chapter8.html#regularisation-and-hyperparameters}{8.12}
of the lecture notes).  Find the optimal results for the hyperparameters
$\lambda$ and the learning rates $\gamma$ and neural network architecture
and compare the $L_2$ results with Ridge regression from project 1 and the
$L_1$ results with the Lasso calculations of project 1.  Use again the same
data sets and the best results from project 1 in your comparisons.

\subsection*{Part e): Classification analysis using neural networks}

With a well-written code it should now be easy to change the activation
function for the output layer.

Here we will change the cost function for our neural network code developed
in parts b)--d) in order to perform a classification analysis (Section
\href{https://educationalmaterialuio.github.io/MachineLearningUiO/doc/LectureNotes/_build/html/chapter8.html#classification-with-neural-networks}{8.11}
of the lecture notes).  The classification problem we will study is the
multiclass MNIST problem of handwritten digits (LeCun, Cortes and Burges),
see the description of the data set at
\url{https://www.tensorflow.org/datasets/catalog/mnist}.  We will use the
softmax cross entropy function discussed in a).  In this project you should
use the \emph{full} MNIST data set ($70\,000$ images of $28\times 28$
pixels), the one used in the Monday lecture of week 41; the code of that
lecture reads exactly the file described below.  Note that the Tuesday
session of week 41 (and Section
\href{https://educationalmaterialuio.github.io/MachineLearningUiO/doc/LectureNotes/_build/html/chapter8.html#examples}{8.14}
of the lecture notes) uses the much smaller digits data set of Scikit-Learn
(\texttt{load\_digits}, $1797$ images of $8\times 8$ pixels, a different data
set of the same kind and not a subset of MNIST).  The reduced set is
convenient for developing and debugging your code, since every run takes
seconds, but the results you report should be for the full MNIST set.

The data set comes with a canonical split into $60\,000$ training images
and $10\,000$ test images, each a $28\times 28$ array of grey-scale pixel
values between 0 and 255.  It can be downloaded, without registering
anywhere, as a single NumPy archive (about 11 MB) from
\url{https://storage.googleapis.com/tensorflow/tf-keras-datasets/mnist.npz}
and loaded with
\begin{lstlisting}
import numpy as np
with np.load("mnist.npz") as data:
    x_train, y_train = data["x_train"], data["y_train"]
    x_test, y_test = data["x_test"], data["y_test"]  # untouched until the end
print(x_train.shape, y_train.shape, x_test.shape, y_test.shape)
# (60000, 28, 28) (60000,) (10000, 28, 28) (10000,)
\end{lstlisting}
A feed-forward network takes a vector as input, so flatten each image to a
vector of $784$ features and scale the pixel values to the interval $[0,1]$
(scaling makes a real difference for gradient descent, see Section 3.13 of
the lecture notes),
\begin{lstlisting}
x_train = x_train.reshape(len(x_train), -1) / 255.0
x_test = x_test.reshape(len(x_test), -1) / 255.0
\end{lstlisting}
Use the canonical split the way it is meant to be used: all training, model
selection and tuning of hyperparameters (learning rate, penalty $\lambda$,
number of layers and nodes, activation functions) is done on the training
data only, for example by setting aside a validation set from the training
data or by cross-validation (Chapter 3 of the lecture notes).  The test set
is used \emph{once}, at the very end, to report the performance of the
models you have selected.  If you are limited by computing time, you may
train on a subset of the training images, but evaluate on the full test
set.

Remember to cite the data set in your report, as you would cite any other
source; the citation is given on the data set page above (Y. LeCun,
C. Cortes and C. J. C. Burges, \emph{The MNIST database of handwritten
digits}, \url{http://yann.lecun.com/exdb/mnist/}).

Feel free to suggest other data sets.  If you find the classic MNIST data
set somewhat limited, try the Fashion-MNIST data set of Zalando Research at
\url{https://github.com/zalandoresearch/fashion-mnist}.  It has exactly the
same format and the same $60\,000/10\,000$ split, with ten classes of
clothing items instead of digits, and the download links, loading code and
the reference to cite are all on that page.

To measure the performance of our classification problem we will use the
so-called \emph{accuracy} score (Section 5.12 of the lecture notes).  The
accuracy is as you would expect just the number of correctly guessed targets
$t_i$ divided by the total number of targets, that is
\[
  \text{Accuracy} = \frac{\sum_{i=1}^n I(t_i = y_i)}{n} ,
\]
where $I$ is the indicator function, $1$ if $t_i = y_i$ and $0$ otherwise.
Here $t_i$ represents the target and $y_i$ the outputs of your FFNN code and
$n$ is simply the number of targets $t_i$.

Discuss your results and give a critical analysis of the various parameters,
including hyperparameters like the learning rates and the regularization
parameter $\lambda$, various activation functions, number of hidden layers
and nodes and activation functions.

Again, we strongly recommend that you compare your own neural network code
for classification and pertinent results against a similar code using
\textbf{Scikit-Learn} or \textbf{TensorFlow/Keras} or \textbf{PyTorch} (see
the paragraph on established libraries in Section
\href{https://educationalmaterialuio.github.io/MachineLearningUiO/doc/LectureNotes/_build/html/chapter8.html#examples}{8.14}
of the lecture notes).

You should also compare your neural network results with those from
logistic regression (Sections 5.3 and 5.6 of the lecture notes).  Logistic
regression is the \emph{baseline} for the classification part: it is the
classifier with no hidden layers, and a network with one or more hidden
layers is only worth its extra cost if you can show that it performs better
than this baseline.  Your neural network code implements the equivalent of
logistic regression by simply setting the number of hidden layers to zero
and keeping just the input and the output layers (with the softmax activation
and the cross entropy cost from part a)).

In addition to the accuracy, you should compute and interpret the
\emph{confusion matrix} of your best classification model (Section 5.12 of
the lecture notes, see also
\url{https://www.researchgate.net/figure/Confusion-matrix-of-MNIST-and-F-MNIST-embeddings_fig5_349758607}).
The confusion matrix shows which digits are confused with which, something a
single accuracy number hides, and if you also report the precision, the
recall and the $F_1$ score, the confusion matrix is what makes these numbers
interpretable.  Scikit-Learn's \texttt{confusion\_matrix} and
\texttt{ConfusionMatrixDisplay} from \texttt{sklearn.metrics} are
convenient for this.

If you wish to compare with logistic regression from
\textbf{Scikit-Learn} as well, the following code uses the flattened and
scaled arrays from above
\begin{lstlisting}
from sklearn.linear_model import LogisticRegression
# Initialize the model (multinomial logistic regression with the saga solver)
model = LogisticRegression(solver='saga', max_iter=1000, random_state=42)
# Train the model
model.fit(x_train, y_train)
from sklearn.metrics import accuracy_score
# Make predictions on the test set
y_pred = model.predict(x_test)
# Calculate accuracy
accuracy = accuracy_score(y_test, y_pred)
print(f"Model Accuracy: {accuracy:.4f}")
\end{lstlisting}

\subsection*{Part f): Critical evaluation of the various algorithms}

After all these glorious calculations, you should now summarize the various
algorithms and come with a critical evaluation of their pros and cons.
Which algorithm works best for the regression case and which is best for
the classification case?  Section
\href{https://educationalmaterialuio.github.io/MachineLearningUiO/doc/LectureNotes/_build/html/chapter8.html#limitations-and-a-top-down-view}{8.15}
of the lecture notes gives a top-down view of what neural networks can and
cannot do that you may find useful here.  These codes can also be part of
your final project 3, but now applied to other data sets.

\section*{Summary of methods to implement and analyze}

\textbf{Required implementation:}
\begin{enumerate}[nosep]
  \item Reuse the regression code and results from project 1; these will act
        as a benchmark for seeing how suited a neural network is for this
        regression task.
  \item Implement a neural network with
  \begin{itemize}[nosep]
    \item a flexible number of layers;
    \item a flexible number of nodes in each layer;
    \item a changeable activation function in each layer (Sigmoid, ReLU,
          Leaky ReLU, as well as linear and softmax);
    \item a changeable cost function, which will be set to the MSE for
          regression and the cross entropy for multiclass classification;
    \item an optional $L_1$ or $L_2$ norm of the weights and biases in the
          cost function (only used for computing gradients, not as
          interpretable metrics).
  \end{itemize}
  \item Implement the back propagation algorithm to compute the gradient of
        your neural network.
  \item Reuse the implementation of plain and stochastic gradient descent
        from project 1 (and adapt the code to work with your neural
        network)
  \begin{itemize}[nosep]
    \item with no optimization algorithm;
    \item with RMSprop;
    \item with Adam.
  \end{itemize}
  \item Implement logistic regression as the baseline classification model
        (equivalent to a neural network with just the output layer, that is
        with zero hidden layers).
  \item Implement scaling and train-test splitting of your data, preferably
        using Scikit-Learn (MNIST comes with its own test set; tune on the
        training data and use the test set only for the final evaluation).
  \item Implement and compute metrics like the MSE, the accuracy and the
        confusion matrix.
\end{enumerate}

\textbf{Required analysis:}
\begin{enumerate}[nosep]
  \item Briefly show and argue for the advantages and disadvantages of the
        methods from project 1.
  \item Explore and show the impact of changing the number of layers, nodes
        per layer, choice of activation function, and inclusion of $L_1$ and
        $L_2$ norms.  Present only the most interesting results from this
        exploration.  2D heat maps will be good for this: start with finding
        a well performing set of hyperparameters, then change two at a time
        in a range that shows good and bad performance.
  \item Show and argue for the advantages and disadvantages of using a
        neural network for regression on your data.
  \item Show and argue for the advantages and disadvantages of using a
        neural network for classification on your data.  In particular, show
        whether and why a network with hidden layers performs better than
        the logistic-regression baseline (no hidden layers).
  \item Compute and interpret a confusion matrix of your best classification
        model (Section 5.12 of the lecture notes and
        \url{https://www.researchgate.net/figure/Confusion-matrix-of-MNIST-and-F-MNIST-embeddings_fig5_349758607}).
  \item Show and argue for the advantages and disadvantages of the different
        gradient methods and learning rates when training the neural
        network.
\end{enumerate}

\section*{Background literature}

\begin{itemize}[nosep]
  \item The lecture notes of the course, in particular Chapter 8 (neural
        networks, at
        \url{https://educationalmaterialuio.github.io/MachineLearningUiO/doc/LectureNotes/_build/html/chapter8.html}),
        Chapter 4 (stochastic gradient descent, RMSprop, Adam and automatic
        differentiation) and Chapter 5 (the cross entropy, the softmax
        function and how to measure the quality of a classifier).
  \item The text of Michael Nielsen is highly recommended, see Nielsen's
        book at \url{http://neuralnetworksanddeeplearning.com/}.  It is an
        excellent read.
  \item Goodfellow, Bengio and Courville, Deep Learning, at
        \url{https://www.deeplearningbook.org/}.  Here we recommend chapters
        6, 7 and 8.
  \item Raschka et al. at
        \url{https://sebastianraschka.com/blog/2022/ml-pytorch-book.html}.
        Here we recommend chapters 11, 12 and 13.
\end{itemize}

\section*{Introduction to numerical projects}

Here follows a brief recipe and recommendation on how to answer the various
questions when preparing your answers.
\begin{itemize}[nosep]
  \item Give a short description of the nature of the problem and the
        eventual numerical methods you have used.
  \item Describe the algorithm you have used and/or developed.  Here you may
        find it convenient to use pseudocoding.  In many cases you can
        describe the algorithm in the program itself.
  \item Include the source code of your program.  Comment your program
        properly.  You should have the code at your GitHub/GitLab link.  You
        can also place the code in an appendix of your report.
  \item If possible, try to find analytic solutions, or known limits in order
        to test your program when developing the code.  In this project the
        gradient from automatic differentiation and the closed-form OLS
        solution of project 1 are exactly such tests for your neural network
        code.
  \item Include your results either in figure form or in a table.  Remember
        to label your results.  All tables and figures should have relevant
        captions and labels on the axes.
  \item Try to evaluate the reliability and numerical stability/precision of
        your results.  If possible, include a qualitative and/or quantitative
        discussion of the numerical stability, eventual loss of precision
        etc.
  \item Try to give an interpretation of your results in your answers to the
        problems.
  \item Critique: if possible include your comments and reflections about the
        exercise, whether you felt you learnt something, ideas for
        improvements and other thoughts you've made when solving the
        exercise.  We wish to keep this course at the interactive level and
        your comments can help us improve it.
  \item Try to establish a practice where you log your work at the computer
        lab.  You may find such a logbook very handy at later stages in your
        work, especially when you don't properly remember what a previous
        test version of your program did.  Here you could also record the
        time spent on solving the exercise, various algorithms you may have
        tested or other topics which you feel worthy of mentioning.
\end{itemize}

\section*{Format for electronic delivery of report and programs}

The preferred format for the report is a PDF file.  You can also use DOC or
postscript formats or an IPython notebook file.  As programming language we
prefer that you choose between C/C++, Fortran2008, Julia or Python.  The
following prescription should be followed when preparing the report:
\begin{itemize}[nosep]
  \item Use Canvas to hand in your projects, log in at
        \url{https://www.uio.no/english/services/it/education/canvas/} with
        your normal UiO username and password.
  \item Upload \textbf{only} the report file or the link to your GitHub/GitLab
        or similar type of repository!  For the source code file(s) you have
        developed please provide us with your link to your GitHub/GitLab or
        similar domain.  The report file should include all of your
        discussions and a list of the codes you have developed.  Do not
        include library files which are available at the course homepage,
        unless you have made specific changes to them.
  \item In your GitHub/GitLab or similar repository, please include a folder
        which contains selected results.  These can be in the form of output
        from your code for a selected set of runs and input parameters.
\end{itemize}
Finally, we encourage you to collaborate.  Optimal working groups consist of
2--3 students.  You can then hand in a common report.

\section*{Software and needed installations}

If you have Python installed (we recommend Python 3) and you feel pretty
familiar with installing different packages, we recommend that you install
the following Python packages via \texttt{pip} as
\begin{quote}\small
\texttt{pip install numpy scipy matplotlib ipython jupyter scikit-learn}\\
\texttt{pip install pandas seaborn}\\
\texttt{pip install jax jaxlib}
\end{quote}
(\texttt{autograd} is a lighter alternative to JAX if you want to check your
back propagation code against automatic differentiation; PyTorch or
TensorFlow are needed only if you choose them for the comparison with an
established library).  For Python 3, replace \texttt{pip} with
\texttt{pip3} where needed.

For OSX users we recommend also, after having installed Xcode, to install
\texttt{brew}.  Brew allows for a seamless installation of additional
software via for example \texttt{brew install python3}.  For Linux users,
with its variety of distributions like for example the widely popular Ubuntu
distribution, you can use \texttt{pip} as well and simply install Python as
\texttt{sudo apt-get install python3} etc.

If you don't want to install various Python packages with their dependencies
separately, we recommend \href{https://docs.anaconda.com/}{Anaconda}, an open
source distribution of the Python and R programming languages for large-scale
data processing, predictive analytics, and scientific computing, that aims to
simplify package management and deployment.  Package versions are managed by
the package management system \texttt{conda}.

Popular software packages written in Python for ML are
\href{http://scikit-learn.org/stable/}{Scikit-learn},
\href{https://jax.readthedocs.io}{JAX},
\href{https://www.tensorflow.org/}{TensorFlow},
\href{http://pytorch.org/}{PyTorch} and \href{https://keras.io/}{Keras}.
These are all freely available at their respective GitHub sites.  They
encompass communities of developers in the thousands or more, and the number
of code developers and contributors keeps increasing.

\end{document}
